\documentclass[10pt]{article}
\usepackage[T1]{fontenc}
\usepackage[utf8]{inputenc}
\usepackage[french]{babel}
\usepackage{amsmath,amssymb,mathrsfs,mathpazo}
\usepackage{geometry}
\geometry{a4paper,margin=2cm}
\usepackage{graphicx,array,tabularx,tabularray,colortbl,multirow,multicol,booktabs,enumitem,pifont,ulem,fancybox,diagbox,cancel}
\usepackage{pstricks,pstricks-add,pst-plot,pst-tree,pst-node,pst-text,pst-func,pst-math,pst-bspline,pst-3dplot,pst-eucl,pst-solides3d}
\let\pscalculate\undefined
\usepackage{pst-bezier}
\usepackage{tikz,pgfplots,tkz-tab}
\usetikzlibrary{arrows,arrows.meta,patterns,calc,positioning,trees,decorations.markings,shapes.geometric,shapes.misc,shapes.arrows,intersections,decorations.pathmorphing,decorations.pathreplacing,backgrounds,shadows,babel}
\pgfplotsset{compat=1.18}
\usepgfplotslibrary{fillbetween}
\usepackage[np]{numprint}
\usepackage{listings,makecell,pas-tableur,twoopt}
\let\barre\undefined
\usepackage[pstricks]{tablvar}
\definecolor{darkgreen}{RGB}{0,100,0}
\providecommand{\red}{\color{red}}
\providecommand{\blue}{\color{blue}}
\providecommand{\green}{\color{green}}
\providecommand{\black}{\color{black}}
\providecommand{\white}{\color{white}}
\providecommand{\gray}{\color{gray}}
\providecommand{\cyan}{\color{cyan}}
\providecommand{\magenta}{\color{magenta}}
\providecommand{\yellow}{\color{yellow}}
\providecommand{\violet}{\color{violet}}
\providecommand{\orange}{\color{orange}}
\providecommand{\euro}{\texteuro}
\providecommand{\EUR}{\texteuro}
\providecommand{\ieme}{\textsuperscript{e}}
\providecommand{\degres}{\ensuremath{^\circ}}
\providecommand{\degree}{\ensuremath{^\circ}}
\providecommand{\vv}[1]{\overrightarrow{#1}}
\providecommand{\Coord}[1]{#1}
\providecommand{\decosix}{}
\providecommand{\starredbullet}{\ensuremath{\bullet}}
\pagestyle{empty}
\setlength{\parindent}{0pt}
\usepackage{adjustbox}\DeclareUnicodeCharacter{2236}{:}\DeclareUnicodeCharacter{2190}{\ensuremath{\leftarrow}}
\providecommand{\R}{\mathbb{R}}
\providecommand{\N}{\mathbb{N}}
\providecommand{\Z}{\mathbb{Z}}
\providecommand{\Q}{\mathbb{Q}}
\providecommand{\D}{\mathbb{D}}
\providecommand{\C}{\mathbb{C}}
\providecommand{\e}{{\mathrm{e}}}
\providecommand{\vect}[1]{\overrightarrow{#1}}
\providecommand{\vectt}[1]{\overrightarrow{\mathrm{#1}}}
\providecommand{\vvt}[1]{\overrightarrow{\mathrm{#1}}}
\providecommand{\np}[2][]{{#2~#1}}
\providecommand{\vv}[1]{{\overrightarrow{#1}}}
\renewcommand{\d}{{\mathrm{d}}}
\providecommand{\degres}{{^\circ}}
\providecommand{\degre}{{^\circ}}
\providecommand{\blue}{\color{blue}}
\providecommand{\red}{\color{red}}
\providecommand{\enskip}{\quad}
\providecommand{\llbracket}{\lbrack\!\lbrack}
\providecommand{\rrbracket}{\rbrack\!\rbrack}
\providecommand{\euro}{{\text{€}}}
\providecommand{\no}{{\text{n°}}}
\providecommand{\interval}[2]{{[#1;#2]}}
\providecommand{\intervFO}[2]{{[#1;#2[}}
\providecommand{\dotfill}{\ldots}

\providecommand{\ds}{\displaystyle}
\providecommand{\textding}[1]{#1}
\providecommand{\Oij}{\(\mathrm{O};\vec{\imath},\vec{\jmath}\)}
\providecommand{\Oijk}{\(\mathrm{O};\vec{\imath},\vec{\jmath},\vec{k}\)}
\providecommand{\Ouv}{\(\mathrm{O};\vec{u},\vec{v}\)}
\providecommand{\intervFF}[2]{\ensuremath{[#1;#2]}}
\renewcommand{\euro}{\ifmmode\text{\texteuro}\else\texteuro\fi}
\DeclareUnicodeCharacter{00B0}{\ifmmode^{\circ}\else\textdegree\fi}
\DeclareUnicodeCharacter{20AC}{\ifmmode\text{\texteuro}\else\texteuro\fi}
% Énoncés : Éducation nationale. Transcription et corrigés : APMEP, Denis Vergès et contributeurs. Découpage et classement : CoopMaths, Sébastien Lozano, Éric Elter et contributeurs. Adaptation de la présentation : Maths Explorer.
% CC BY-SA : https://coopmaths.fr/www/about/
% https://coopmaths.fr/alea/static/bac/2026/tex/bac_2026_05_sujet2_ameriquenord_2_cor.tex
\begin{document}
\noindent\textit{APMEP / CoopMaths. CC BY-SA. Présentation adaptée par Maths Explorer.}\par\medskip
\textbf{Partie A: étude du sens de variation d'une fonction}

\medskip

On considère la fonction $f$ définie sur $\R$ par 
$f(x) = \dfrac{2x}{\sqrt{1 + x^2}}$.

\begin{enumerate}
\item On résout l'équation $f(x) = x$.

$\begin{aligned}
f(x)=x
& \iff  \dfrac{2x}{\sqrt{1 + x^2}} =x
\iff x\left ( \dfrac{2}{\sqrt{1 + x^2}} - 1\right )=0
\iff x \left (\dfrac{2-\sqrt{1+x^2}}{\sqrt{1+x^2}}\right )=0
\end{aligned}$

Soit $x=0$, soit $\dfrac{2-\sqrt{1+x^2}}{\sqrt{1+x^2}}=0$.

$\begin{aligned}
\dfrac{2-\sqrt{1+x^2}}{\sqrt{1+x^2}}=0
& \iff 2-\sqrt{1+x^2}=0
\iff 2=\sqrt{1+x^2}
\iff 4=1+x^2
\iff 3=x^2\\
& \iff x=\sqrt{3}  \text{ ou } x=-\sqrt{3}
\end{aligned}$

L'équation $f(x)=x$ a pour ensemble solution $S=\left \{ -\sqrt{3}\,;\,0\,;\,\sqrt{3}\strut \right \}$.

\item
	\begin{enumerate}
		\item On admet que la fonction $f$ est dérivable sur $\R$.

La dérivée de la fonction $x \longmapsto \sqrt{1+x^2}$ est la fonction $x \longmapsto \dfrac{2x}{2\sqrt{1+x^2}}$, soit $x \longmapsto \dfrac{x}{\sqrt{1+x^2}}$.

$\begin{aligned}
f'(x)
& = \dfrac{2\sqrt{1+x^2} - 2x\times\frac{x}{\sqrt{1+x^2}}}{1+x^2}
= \dfrac{\frac{2\left ( \sqrt{1+x^2}\right )^2-2x^2}{\sqrt{1+x^2}}}{1+x^2}
= \dfrac{2+2x^2-2x^2}{\left (1+x^2\right ) \sqrt{1+x^2}}
= \dfrac{2}{\left (1+x^2\right ) \sqrt{1+x^2}}
\end{aligned}$


		\item 
Pour tout réel $x$, $1+x^2>0$ et $\sqrt{1+x^2}>0$, donc $f'(x)>0$.

Donc la fonction $f$ est strictement croissante sur $\R$.
	\end{enumerate}
\end{enumerate}

\medskip

\textbf{Partie B : étude de la convergence d'une suite récurrente}

\medskip

La suite $(u_n)$ est définie par $u_0 = 1$ et, pour tout entier naturel $n,\: u_{n+1} = f(u_n)$.

\medskip

\begin{enumerate}
\item On veut démontrer par récurrence que, pour tout   $n$, la propriété $1 \leqslant u_n \leqslant u_{n+1} < \sqrt 3$ est vraie.

\begin{list}{\textbullet}{}
\item \textbf{Initialisation}

$u_0=1$ et $u_1=f(u_0)=f(1)=\dfrac{2\times 1}{\sqrt{1+1^2}}=\dfrac{2}{\sqrt{2}}=\sqrt{2}<\sqrt{3}$

Donc la propriété  $1 \leqslant u_0 \leqslant u_1 < \sqrt{3}$ est vraie pour $n=0$.

\item \textbf{Hérédité}

On suppose que la propriété est vraie au rang $n$, c'est-à-dire: $1 \leqslant u_n \leqslant u_{n+1} < \sqrt 3$.

La fonction $f$ est strictement croissante sur $\R$ donc:
$f(1) \leqslant f\left (u_n\right ) \leqslant f\left (u_{n+1}\right ) < f\left (\sqrt 3\right )$

$f(1) = \sqrt{2}>1$;
$f\left (u_n\right )=u_{n+1}$;
$f\left (u_{n+1}\right )=u_{n+2}$ et

$f\left (\sqrt{3}\right ) = \dfrac{2\times \sqrt{3}}{\sqrt{1+\left (\sqrt{3}\right )^2}}=\dfrac{2\sqrt{3}}{\sqrt{4}}=\sqrt{3}$

Donc on a: $1 \leqslant u_{n+1} \leqslant u_{n+2} < \sqrt 3$, et donc la propriété est vraie au rang $n+1$.

\item \textbf{Conclusion}

La propriété est vraie au rang 0, et elle est héréditaire pour tout $n\geqslant 0$. Elle est donc vraie pour tout entier naturel $n$.
\end{list}

On a donc démontré que, pour tout $n$, on a: $1 \leqslant u_n \leqslant u_{n+1} < \sqrt 3$.

\item 
\begin{list}{\textbullet}{}
\item Pour tout $n$, on a: $u_n \leqslant u_{n+1}$ donc la suite $(u_n)$ est croissante.
\item Pour tout $n$, on a: $u_n <  \sqrt{3}$ donc la suite $(u_n)$ est majorée.
\end{list}

La suite $(u_n)$ est croissante et majorée donc, d'après le théorème de la convergence monotone, on peut déduire que la suite $(u_n)$ est convergente.

On appelle $\ell$ la limite de $(u_n)$. Comme pour tout $n$, $1 \leqslant u_n < \sqrt{3}$, on a: $1 \leqslant \ell \leqslant \sqrt{3}$.

\begin{list}{\textbullet}{}
\item $\ds\lim_{n\to +\infty} u_n=\ell$ et on sait que la fonction $f$ est dérivable sur $\R$ donc continue; on en déduit que $\ds\lim_{n\to +\infty}f(u_n) = f(\ell)$. Or $u_{n+1}=f(u_n)$ donc $\ds\lim_{n\to +\infty} u_{n+1}=f(\ell)$.

\item La limite d'une suite est unique, donc $\ds\lim_{n\to +\infty}u_n = \ds\lim_{n\to +\infty} u_{n+1}$. 
\end{list}

On en déduit que $\ell=f(\ell)$ et  donc que  $\ell$ est solution de l'équation $f(x)=x$.





La seule solution vérifiant $1 \leqslant \ell \leqslant \sqrt{3}$ est $\sqrt{3}$ (voir \textbf{A.1}), donc $\ds\lim_{n\to +\infty} u_n=\sqrt{3}$.

\item 
Pour tout entier naturel $n$, on pose :
$v_n = \dfrac{u_n^2}{3 - u_n^2}$.

On admet que la suite $(v_n)$ est bien définie.
	\begin{enumerate}
		\item 
\begin{list}{\textbullet}{}
\item $u_0=1$ donc $v_0=\dfrac{u_0^2}{3-u_0^2} = \dfrac{1^2}{3-1^2}=\dfrac{1}{2}=0,5$
\item
$\begin{aligned}[t]
v_{n+1}
& = \dfrac{u_{n+1}^2}{3-u_{n+1}^2}
= \dfrac{\left ( \frac{2u_n}{\sqrt{1+u_n^2}}\right )^2}{3-\left ( \frac{2u_n}{\sqrt{1+u_n^2}}\right )^2}
= \dfrac{ \frac{4u_n^2}{1+u_n^2}}{3-\frac{4u_n^2}{1+u_n^2}}
= \dfrac{ \frac{4u_n^2}{1+u_n^2}}{\frac{3\left ( 1+u_n^2\right ) -4u_n^2}{1+u_n^2}}\\
& = \dfrac{4u_n^2}{1+u_n^2}\times \dfrac{1+u_n^2}{3 + 3u_n^2 -4u_n^2}
= \dfrac{4u_n^2}{3-u_n^2}
= 4\times \dfrac{u_n^2}{3-u_n^2}
= 4v_n
\end{aligned}$
\end{list}

Donc la suite $(v_n)$ est géométrique de premier terme $v_0=0,5$ et de raison $q=4$.

		\item
\begin{list}{\textbullet}{}
\item La suite $(v_n)$ est géométrique de premier terme $v_0=0,5$ et de raison $q=4$ donc, pour tout $n$, on a: $v_n=v_0\times q^n = 0,5 \times 4^n$.

On en déduit que, pour tout $n$, $v_n>0$.

\item 
$\begin{aligned}[t]
v_n = \dfrac{u_n^2}{3 - u_n^2}
& \iff v_n \left (3-u_n^2\right ) = u_n^2
\iff 3v_n -v_n\times u_n^2 = u_n^2
\iff 3v_n = u_n^2\left (1+v_n\right )\\
& \iff \dfrac{3v_n}{1+v_n} =u_n^2
\end{aligned}$

$v_n>0$ donc $1+v_n>0$ donc $\dfrac{3v_n}{1+v_n}>0$ et donc $\sqrt{\dfrac{3v_n}{1+v_n}}$ existe.

De plus on sait que $1 \leqslant u_n < \sqrt{3}$ donc $u_n>0$.

On a donc: $u_n=\sqrt{\dfrac{3v_n}{1+v_n}}$.
\end{list}

On en déduit que pour tout $n$, on a:
$u_n=\sqrt{\dfrac{3\times 0,5 \times 4^n}{1+0,5\times 4^n}}
=\sqrt{\dfrac{1,5 \times 4^n}{1+0,5\times 4^n}}$

		\item
$u_n = \sqrt{\dfrac{1,5 \times 4^n}{1+0,5\times 4^n}} = \sqrt{\dfrac{1,5 \times 4^n}{4^n \left ( \frac{1}{4^n}+0,5\right )}} = \sqrt{\dfrac{1,5}{ \frac{1}{4^n}+0,5}}$		

D'après les propriétés des suites géométriques, on sait que si $q>1$, alors $\ds\lim_{n\to +\infty} q^n=+\infty$.

4>1 donc $\ds\lim_{n\to +\infty}4^n=+\infty$ donc
$\ds\lim_{n\to +\infty} \dfrac{1}{4^n} =0$ donc
$\ds\lim_{n\to +\infty} \dfrac{1,5}{\frac{1}{4^n}+0,5}= \dfrac{1,5}{0,5} = 3$

Par continuité de la fonction racine carrée, on en déduit que: $\ds\lim_{n\to +\infty} u_n=\sqrt{3}$.
	\end{enumerate}
\end{enumerate}

\medskip

\textbf{Partie C : étude de la convergence de la somme de termes}

\medskip

Pour tout entier naturel $n$ non nul, on pose $S_n = u_0^2 + u_1^2 + \ldots + u_{n-1}^2$.

\medskip

\begin{enumerate}
\item On complète le script Python ci-dessous afin que celui-ci permette de lister les $p$ premiers termes de la suite $(S_n)$.

\begin{center}
\begin{ttfamily}
\begin{tabular}{|l|}\hline
\textbf{from} math \textbf{import}*\\
~\\
\textbf{def} termes(p) :\\
\qquad  u = \blue 1\\
\qquad  S = 0\\
\qquad  L = [ ]\\
\qquad  \textbf{for} i \textbf{in} range(p)\\
\qquad \qquad S = \blue{}S+u**2\\
\qquad \qquad u = \blue{}2*u/sqrt(1+u**2)\\
\qquad \qquad L.append(S) \\
\qquad \textbf{return} L\\ \hline
\end{tabular}
\end{ttfamily}
\end{center}








\item On rappelle que, pour tout entier naturel $k$, on a $1 \leqslant u_k \leqslant \sqrt 3$. Donc: $1 \leqslant u_k^2 \leqslant 3$.

$S_n$ est la somme de $n$ termes compris entre 1 et 3; donc $n\times 1 \leqslant S_n \leqslant n\times 3$ donc $n \leqslant S_n \leqslant 3n$.



\item
\begin{list}{\textbullet}{Calculs de limites}
\item $n \leqslant S_n$ et $\ds\lim_{n\to +\infty} n=+\infty$ donc, par comparaison, $\ds\lim_{n\to +\infty} S_n=+\infty$.
\item $n \leqslant S_n \leqslant 3n$ donc $\dfrac{n}{n^2} \leqslant \dfrac{S_n}{n^2} \leqslant \dfrac{3n}{n^2}$  donc $\dfrac{1}{n} \leqslant \dfrac{S_n}{n^2} \leqslant \dfrac{3}{n}$.

$\ds\lim_{n\to +\infty} \dfrac{1}{n}=0$ et $\ds\lim_{n\to +\infty} \dfrac{3}{n}=0$ donc, d'après le théorème des gendarmes, $\ds\lim_{n\to +\infty} \dfrac{S_n}{n^2}=0$.
\end{list}

\end{enumerate}

\medskip

\end{document}
